\documentclass[12pt]{article}
\usepackage[utf8]{inputenc}
\usepackage[spanish, es-noshorthands]{babel}
\usepackage{csquotes} % para que no se queje babel
\decimalpoint
\usepackage{amssymb,amsmath,amsthm,amsfonts,mathtools}
\usepackage{tcolorbox}
\usepackage{microtype}

% ---- LINKS ----
\usepackage{hyperref}
\hypersetup{colorlinks=true, urlcolor=blue, linkcolor=VerdeITAM, citecolor=VerdeITAM}

% ---- IMÁGENES ----
\usepackage{graphicx}
\graphicspath{ {imgs/} }

% ---- FORMATO DE TÍTULOS ----
\usepackage{titlesec}
\titlespacing{\section}{0pt}{1em plus 4pt minus 2pt}{0pt plus 2pt minus 2pt}
\titlespacing{\subsection}{0pt}{1em plus 4pt minus 2pt}{0pt plus 2pt minus 2pt}
\titlespacing{\subsubsection}{0pt}{1em plus 4pt minus 2pt}{0pt plus 2pt minus 2pt}

% ---- FORMATO DE LISTAS ----
\usepackage{enumitem}
\setlist[itemize]{topsep=-0.5em} % es media m para compensar
\setlist[enumerate]{topsep=-0.5em}

% ---- GRÁFICAS ----
\usepackage{tikz}
\usepackage{scalerel}
\usepackage{pict2e}
\usepackage{tkz-euclide}
\usetikzlibrary{calc}
\usetikzlibrary{patterns,arrows.meta}                   % Mejores flechas
\usetikzlibrary{shadows}                                % Sombras
\usetikzlibrary{external}                               % ¿?
\usetikzlibrary{decorations.pathreplacing,calligraphy}  % Para poder graficar el '{'
% Plots
\usepackage{pgfplots}
\pgfplotsset{compat=newest}                             % Ancho de las gráficas: width=6.5cm,
\usepgfplotslibrary{statistics}
\usepgfplotslibrary{fillbetween}
% ---- PORTADA ----
\usetikzlibrary{ shapes.geometric }
\usepackage{anyfontsize}

% ---- COLORES ----
\usepackage{xcolor}
\definecolor{miAzul}{RGB}{13, 33, 161}
\definecolor{miRojo}{RGB}{186, 24, 27}
\definecolor{VerdeITAM}{RGB}{4, 107, 83}

% ---- FORMATO DE LAS PÁGINAS ----
\usepackage[a4paper, width=165mm, top=20mm, bottom = 20mm]{geometry}
\usepackage{fancyhdr}
\pagestyle{fancy}
% Primera página. Comentar si se quiere que todas sean iguales
% \fancypagestyle{firststyle}{
%    \fancyhf{}
%    \fancyfoot[R]{\footnotesize \thepage}
%    \renewcommand{\headrulewidth}{0pt}
% }
% Para el resto del documento.
\fancyhf{}
% \lhead{\footnotesize Subtítulo}
\lhead{\footnotesize \textbf{Funciones trigonométricas en Ensamblador x86}}
\rhead{\footnotesize \today}
\lfoot{\footnotesize J. Alberto Márquez L.}
\cfoot{\footnotesize \textsc{Organización y Programación de Computadoras}}
\rfoot{\footnotesize \thepage}

\renewcommand{\headrulewidth}{0.4pt}
\renewcommand{\footrulewidth}{0.4pt}
\setlength{\headheight}{15pt}

\setlength{\parskip}{1em}
\setlength{\parindent}{0cm}

% ---- TABLAS -----
\usepackage{tabularx}

% ---- GRADOS -----
\usepackage{textcomp}
\usepackage{gensymb}

% ---- BIBLIOGRAFÍA ----
\usepackage[backend=biber, style=apa]{biblatex}
\addbibresource{referencias.bib}

\title{
  Organización y Programación de Computadoras \\
  \large Funciones trigonométricas en Ensamblador x86
}
\author{José Alberto Márquez Luján, 187917}
\date{\today}

\begin{document}
\thispagestyle{empty}

\begin{tikzpicture}[remember picture,overlay]
  %%%%%%%%%%%%%%%%%%%% Background %%%%%%%%%%%%%%%%%%%%%%%%
  \fill[VerdeITAM] (current page.south west) rectangle (current page.north
  east);

  \foreach \i in {2.5,...,22} { \node[rounded corners,VerdeITAM!60,draw,regular
    polygon,regular polygon sides=6, minimum size=\i cm,ultra thick] at
  ($(current page.west)+(2.5,-5)$) {} ; }

  %%%%%%%%%%%%%%%%%%%% Background Polygon %%%%%%%%%%%%%%%%%%%%
  \foreach \i in {0.5,...,22} { \node[rounded corners,VerdeITAM!60,draw,regular
    polygon,regular polygon sides=6, minimum size=\i cm,ultra thick] at
  ($(current page.north west)+(2.5,0)$) {} ; }

  \foreach \i in {0.5,...,22} { \node[rounded corners,VerdeITAM!90,draw,regular
    polygon,regular polygon sides=6, minimum size=\i cm,ultra thick] at
  ($(current page.north east)+(0,-9.5)$) {} ; }

  \foreach \i in {21,...,6} { \node[VerdeITAM!85,rounded corners,draw,regular
    polygon,regular polygon sides=6, minimum size=\i cm,ultra thick] at
  ($(current page.south east)+(-0.2,-0.45)$) {} ; }

  %%%%%%%%%%%%%%%%%%%% Title of the Report %%%%%%%%%%%%%%%%%%%%
  \node[
    left, VerdeITAM!5,
    minimum width=0.625*\paperwidth,
    minimum height=3cm,
    text width=0.625*\paperwidth,
    align=flush right,
    rounded corners
  ] at ($(current page.north east)+(-0.5,-9.5)$) {
    {\large \textsc{Organización y Programación \\ de Computadoras}}
  };

  %%%%%%%%%%%%%%%%%%%% Subtitle %%%%%%%%%%%%%%%%%%%%
  \node[
    left, VerdeITAM!10,
    minimum width=0.625*\paperwidth,
    minimum height=2cm,
    text width=0.625*\paperwidth,
    align=flush right,
    rounded corners
  ] at ($(current page.north east)+(-0.5,-11)$) {
    {\fontsize{25}{30} \selectfont \bfseries Funciones trigonométricas \\ en Ensamblador x86}
  };

  %%%%%%%%%%%%%%%%%%%% Author Name %%%%%%%%%%%%%%%%%%%%
  \node[
    left, VerdeITAM!5,
    minimum width=0.625*\paperwidth,
    minimum height=2cm,
    text width=0.625*\paperwidth,
    align=flush right,
    rounded corners
  ] at ($(current page.north east)+(-0.5,-13)$) {
    {\Large José Alberto Márquez Luján, 187917}
  };

  %%%%%%%%%%%%%%%%%%%% Year %%%%%%%%%%%%%%%%%%%%
  \node[rounded corners,fill=VerdeITAM!70,text=VerdeITAM!5,regular
  polygon,regular polygon sides=6, minimum size=2.5 cm,inner sep=0,ultra thick]
  at ($(current page.west)+(2.5,-5)$) {\LARGE \bfseries 2026};

\end{tikzpicture}

% \maketitle
% \thispagestyle{firststyle}

\newpage
\section{Introducción}

Este reporte describe el diseño, implementación y análisis de un programa
escrito enteramente en lenguaje ensamblador x86 que calcula las funciones
trigonométricas seno y coseno. Los cálculos se hacen con la expansión de series
de Taylor y se evalúan en la unidad de aritmética de punto flotante (FPU por
sus siglas en inglés: \textit{Floating-Point Unit}). Para conocimientos sobre
el lenguaje ensamblador x86, seguimos a \textcite{IrvineBook} y las notas de la
clase de Organización y programación de computadoras.

\section{Teoría}

\subsection{Lenguaje ensamblador}

\subsection{Tipos de datos enteros}

\subsection{Tipos de datos de punto flotante}

\subsection{Arquitectura del FPU y registros}

\section{Conceptos matemáticos}

\subsection{Series de Taylor para seno y coseno}

Según \textcite{ThomasBook}, si $f$ es una función con derivadas de todos los
órdenes en algún intervalo que contenga a $a$ como punto interior, entonces la
\textbf{serie de Taylor generada por $f$ en $x=a$} es
\begin{equation}
  \sum_{k=0}^{\infty} \frac{f^{(k)}(a)}{k!} (x-a)^k = f(a) + f'(a)(x-a) +
    \frac{f''(a)}{2!}(x-a)^2 + \cdots.
\end{equation}
y la \textbf{serie de Maclaurin generada por $f$} es
\begin{equation}
  \sum_{k=0}^{\infty} \frac{f^{(k)}(0)}{k!} x^k = f(0) + f'(0)x +
    \frac{f''(0)}{2!}x^2 + \cdots,
\end{equation}
que es la serie de Taylor generada por $f$ en $x=0$.

\textcite{ThomasBook} menciona que, con frecuencia, la serie de Maclaurin
generada por $f$ sólo se llama serie de Taylor de $f$. Nosotros seguiremos esa
convención en este reporte.

No es difícil ver que la serie de Taylor para $\sin(x)$ en $x=0$ es
\begin{equation}
  \sum_{k=0}^{\infty} = \frac{(-1)^k x^{2k+1}}{(2k+1)!}
    = x - \frac{x^3}{3!} + \frac{x^5}{5!} - \cdots,
\end{equation}
mientras que la serie de Taylor para $\cos(x)$ en $x=0$ es
\begin{equation}
  \sum_{k=0}^{\infty} = \frac{(-1)^k x^{2k}}{(2k)!}
    = 1 - \frac{x^2}{2!} + \frac{x^4}{4!} - \cdots.
\end{equation}

\subsection{Residuos y convergencia}

El teorema 24 de \textcite{ThomasBook} nos dice que, si existe una constante
positiva $M$ tal que $\vert f^{(n+1)}(t) \vert \leq M$ para toda $t$ entre $x$
y $a$, inclusive, entonces el término residual $R_n(x)$ del teorema de Taylor
satisface la desigualdad
\begin{equation} \label{residuo}
  \vert R_n(x) \vert \leq M \frac{ \vert x - a \vert^{n+1} }{ (n+1)! }.
\end{equation}
y, si esta desigualdad se cumple para toda $n$ y $f$ satisface todas las demás
condiciones del teorema de Taylor, entonces la serie converge a $f(x)$.

Como todas las derivadas de $\sin(x)$ tienen valores absolutos menores o
iguales que 1, entonces podemos aplicar \eqref{residuo} con $M = 1$ para
obtener
\begin{equation} \label{residuo-sin}
  \vert R_{2k+1}(x) \vert \leq \frac{ \vert x \vert^{2k+2} }{ (2k+2)! }.
\end{equation}
Como el lado derecho de \eqref{residuo-sin} tiende a cero cuando $k \to
\infty$, entonces $\vert R_{2k+1}(x) \to 0$ y la serie de Taylor para $\sin(x)$
converge a $\sin(x)$ para toda $x$. Por lo tanto,
\begin{equation}
  \sin(x) = \sum_{k=0}^{\infty} = \frac{(-1)^k x^{2k+1}}{(2k+1)!}
  = x - \frac{x^3}{3!} + \frac{x^5}{5!} - \cdots.
\end{equation}

Análogamente, podemos aplicar \eqref{residuo} con $M = 1$ para $\cos(x)$:
\begin{equation} \label{residuo-cos}
  \vert R_{2k}(x) \vert \leq \frac{ \vert x \vert^{2k+1} }{ (2k+1)! }.
\end{equation}
De la misma manera que antes, $R_{2k} \to 0$ cuando $k \to \infty$. Por lo
tanto, la serie converge a $\cos(x)$ para todo valor de $x$, lo que quiere
decir que
\begin{equation}
  \cos(x) = \sum_{k=0}^{\infty} = \frac{(-1)^k x^{2k}}{(2k)!}
  = 1 - \frac{x^2}{2!} + \frac{x^4}{4!} - \cdots.
\end{equation}

Notemos que, aunque ambas series convergen para toda $x$ ---lo que quiere decir
que eventualmente producirán un resultado cercano a la función, sin importar la
entrada---, la convergencia depende de manera importante de $\vert x \vert$. Si
$\vert x \vert < 1$, cada término sucesivo es dramáticamente más pequeño que el
anterior, por lo que las series convergen rápidamente. Aunque no es necesario
que $\vert x \vert < 1$, sí conviene que $\vert x \vert$ sea pequeño. Por
ejemplo, para $\sin(x)$ con $k=10$ y $\vert x \vert \leq \frac{\pi}{2}$, por
\eqref{residuo-sin} tenemos que
\begin{equation}
  \vert R_{2(10)+1}(x) \vert \leq \frac{ \vert x \vert^{22} }{ (22)! }
  \leq \frac{ \vert \pi/2 \vert^{22} }{ (22)! }
  \approx 1.836 \times 10^{-17},
\end{equation}
lo cual es una precisión bastante razonable. Sin embargo, si realizamos el
mismo cálculo con $x = 2\pi$, o bien, 360$\degree$, tenemos que
\begin{equation}
  \vert R_{2(10)+1}(2\pi) \vert \leq \frac{ \vert 2\pi \vert^{22} }{ (22)! }
  \approx 3.23 \times 10^{-4},
\end{equation}
lo que implica una pérdida de precisión bastante considerable. Este hecho
motivó, como veremos más adelante, el ajuste de rango al momento de convertir
grados a radianes.

\section{Diseño del programa}

\subsection{Estructura modular}

El código está dividido en 11 archivos con terminación \texttt{.asm}. Cada
archivo contiene uno o dos procedimientos relacionados entre sí. A
continuación, presentamos una tabla con los procedimientos que tiene cada
archivo y una breve descripción.

\begin{center}
  \begin{tabularx}{\textwidth}{p{10em} p{10em} X}
    Módulo & Procedimiento(s) & Descripción \\
    \hline \hline
    \texttt{main.asm} & \texttt{main}, \texttt{ImprimeMenu} &
      Punto de entrada, bucle de menú, texto compartido. \\
    \hline
    \texttt{pide-radianes.asm} & \texttt{PideRadianes} &
      Se encarga de pedir los radianes al usuario y leer la entrada. \\
    \hline
    \texttt{pide-grados.asm} & \texttt{PideGrados} &
      Se encarga de pedir los grados al usuario y leer la entrada. \\
    \hline
    \texttt{tabla-grados.asm} & \texttt{TablaGrados} &
      Se encarga de llamar a las funciones para calcular una tabla con valores de 0$\degree$ a 360$\degree$ en pasos de 10$\degree$. \\
    \hline
    \texttt{eval-sin-cos.asm} & \texttt{EvalSinCos} &
      Dados los radianes, se encarga de llamar a los procedimientos para calcular seno, coseno e identidad trigonométrica. \\
    \hline
    \texttt{eval-grados.asm} & \texttt{EvalGrados} &
      Dados los grados, se encarga de llamar a los procedimientos para convertir a radianes y luego evaluar. \\
    \hline
    \texttt{convierte-rango.asm} & \texttt{ConvierteRangoSin},\newline\texttt{ConvierteRangoCos} &
      Convierte grados a radianes y aplica una reducción de rango. \\
    \hline
    \texttt{imprime-grados.asm} &  \texttt{ImprimeGrados} &
      Se encarga del formato para imprimir los resultados cuando se manejan grados. \\
    \hline
    \texttt{sine-n.asm} &  \texttt{SinFloat} &
      Evalúa $n$ términos de la serie de Taylor para $\sin(x)$, con $x$ en radianes. \\
    \hline
    \texttt{cosine-n.asm} &  \texttt{CosFloat} &
      Evalúa $n$ términos de la serie de Taylor para $\cos(x)$, con $x$ en radianes. \\
    \hline
    \texttt{square-float.asm} &  \texttt{SquareFloat} &
      Se encarga de elevar al cuadrado el valor en ST(0). \\
  \end{tabularx}
\end{center}

\subsection{Diagrama de dependencias}

\begin{figure}[H]
\centering
\resizebox{\textwidth}{!}{
\begin{tikzpicture}[
  node distance=1.4cm and 2.2cm,
  every node/.style={
    draw, rounded corners, align=center,
    font=\small, minimum width=3.8cm, minimum height=1cm,
    fill=white, drop shadow
  },
  edge/.style={->, >=stealth, thick}
]

%% Nodes
\node (main)               {\textbf{main}\\\texttt{main.asm}};

\node (ImprimeMenu)  [below left=of main, xshift=-1cm]   {\textbf{ImprimeMenu}\\\texttt{main.asm}};
\node (PideRadianes) [below=of main, xshift=-2.5cm]        {\textbf{PideRadianes}\\\texttt{pide-radianes.asm}};
\node (PideGrados)   [below=of main, xshift=2.5cm]       {\textbf{PideGrados}\\\texttt{pide-grados.asm}};
\node (TablaGrados)  [below right=of main, xshift=1cm]   {\textbf{TablaGrados}\\\texttt{tabla-grados.asm}};

\node (EvalSinCos)   [below=of PideRadianes] {\textbf{EvalSinCos}\\\texttt{eval-sin-cos.asm}};
\node (EvalGrados)   [below=of PideGrados] {\textbf{EvalGrados}\\\texttt{eval-grados.asm}};
\node (ImprimeGrados)[below=of TablaGrados] {\textbf{ImprimeGrados}\\\texttt{imprime-grados.asm}};

\node (ConvierteRangoSin) [below =of EvalGrados] {\textbf{ConvierteRangoSin}\\\texttt{convierte-rango.asm}};
\node (ConvierteRangoCos) [below =of ImprimeGrados] {\textbf{ConvierteRangoCos}\\\texttt{convierte-rango.asm}};

\node (SinFloat)    [below =of EvalSinCos] {\textbf{SinFloat}\\\texttt{sine-n.asm}};
\node (CosFloat)    [below left=of EvalSinCos] {\textbf{CosFloat}\\\texttt{cosine-n.asm}};
\node (SquareFloat) [below =of CosFloat] {\textbf{SquareFloat}\\\texttt{square-float.asm}};

%% Edges
\draw[edge] (main) -- (ImprimeMenu);
\draw[edge] (main) -- (PideRadianes);
\draw[edge] (main) -- (PideGrados);
\draw[edge] (main) -- (TablaGrados);

\draw[edge] (PideRadianes) -- (EvalSinCos);
\draw[edge] (PideGrados)   -- (EvalGrados);
\draw[edge] (PideGrados)   -- (ImprimeGrados);
\draw[edge] (TablaGrados)  -- (EvalGrados);
\draw[edge] (TablaGrados)  -- (ImprimeGrados);

\draw[edge] (EvalGrados) -- (ConvierteRangoSin);
\draw[edge] (EvalGrados) -- (ConvierteRangoCos);
\draw[edge] (EvalGrados) -- (EvalSinCos);

\draw[edge] (EvalSinCos) -- (SinFloat);
\draw[edge] (EvalSinCos) -- (CosFloat);
\draw[edge] (EvalSinCos) -- (SquareFloat);

\draw[edge] (SinFloat)  -- (SquareFloat);
\draw[edge] (CosFloat)  -- (SquareFloat);

\end{tikzpicture}
}
\caption{Diagrama de dependencias.}
\label{fig:dependencias}
\end{figure}


\subsection{Estrategia de paso de parámetros}

\subsubsection{Valores de punto flotante}

\subsubsection{Valores enteros}

\section{Implementación}

\subsection{Recurrencia en las series de Taylor}

\subsection{Conversión de grados a radianes y reducción de rango}

\subsection{Retos de implementación}

\subsection{Consideraciones}

\section{Capturas de pantalla}

\section{Consideraciones de Linux}

\printbibliography

\end{document}
